\clearpage
\SourceLeaf{018}\pdfbookmark[0]{Contents.}{contenta}
\begin{center}
{\Large CONTENTS.}\par\bigskip
{\large BOOK ONE.}\par\medskip
General relations among the quantities that determine the motions of celestial bodies around the Sun.
\end{center}
\noindent\begin{tabular}{@{}p{.18\linewidth}p{.78\linewidth}@{}}
Section I. & \emph{Relations concerning a single position in the orbit.}\\[.4em]
Section II. & \emph{Relations concerning a single position in space.}\\[.4em]
Section III. & \emph{Relations among several positions in the orbit.}\\[.4em]
Section IV. & \emph{Relations among several positions in space.}
\end{tabular}
\begin{center}\bigskip
{\large BOOK TWO.}\par\medskip
Investigation of the orbits of celestial bodies from geocentric observations.
\end{center}
\noindent\begin{tabular}{@{}p{.18\linewidth}p{.78\linewidth}@{}}
Section I. & \emph{Determination of the orbit from three complete observations.}\\[.4em]
Section II. & \emph{Determination of the orbit from four observations, of which only two are complete.}\\[.4em]
Section III. & \emph{Determination of the orbit that fits any number of observations as closely as possible.}\\[.4em]
Section IV. & \emph{On determining orbits while taking perturbations into account.}
\end{tabular}
\par\bigskip\noindent\textsc{Tables.}
